\section{Taller resuelto: Medir el error de planeo tendido}

\noindent\textbf{Nivel:} Intermedio.

\apartado{1. Enunciado}{Un planeador de 455 kg, S=16.2 m² y polar 0.014+0.032CL² vuela a CL=0.78 en rho=1.05. Compare TAS y descenso exactos con la aproximación L=W y seno de delta igual a tangente.}

\apartado{2. Datos}{Equilibrio estacionario sin empuje; gamma negativo.}

\apartado{3. Incógnitas}{Las magnitudes y comprobaciones solicitadas en el caso.}

\apartado{4. Sistema de referencia}{Trayectoria aérea en calma y fuerzas viento.}

\apartado{5. Hipótesis físicas}{La misma polar se utiliza en ambas comparaciones; sólo cambian las aproximaciones trigonométricas.}

\apartado{6. Modelo matemático}{Medir el error de planeo tendido}

\apartado{7. Ecuaciones simbólicas}{Se parte de los balances y definiciones siguientes.}

\begin{equation}\delta=\arctan(C_D/C_L)\end{equation}

\begin{equation}V_{exacto}=\sqrt{2W\cos\delta/(\rho SC_L)}\end{equation}

\begin{equation}s_{exacto}=V_{exacto}\sin\delta\end{equation}

\begin{equation}V_{aprox}=\sqrt{2W/(\rho SC_L)},\quad s_{aprox}=V_{aprox}C_D/C_L\end{equation}

\apartado{8. Desarrollo algebraico}{Las operaciones se realizan en este orden, conservando las condiciones del modelo.}

\begin{enumerate}

\item Se evalúa un único CD y ángulo para ambas vías.

\item Se conserva coseno en el equilibrio exacto.

\item Se calcula el error relativo de descenso y se interpreta su tamaño respecto a datos de polar.

\end{enumerate}

\apartado{9. Sustitución numérica}{Se conservan los resultados intermedios sin redondear y se muestran todas las operaciones escalares disponibles.}

\textbf{Desarrollo de las magnitudes intermedias.} De $D=W\sin\delta$ y $L=W\cos\delta$, dividir da $\tan\delta=C_D/C_L$. La sustentación exacta usa el coseno; el descenso usa el seno. La aproximación tendida sustituye ambos factores por sus límites de ángulo pequeño.\begin{enumerate}

\item Calcular la polar al CL impuesto.

\[\fitmath{C_D=0.014+0.032\cdot {0.78}^{2}\approx 0.0334688\,}\]

\item Obtener ángulo de descenso positivo.

\[\fitmath{\delta=\arctan\left(\frac{C_D}{0.78}\right)}\]

\[\fitmath{\delta=\arctan\left(\frac{(0.0334688)}{0.78}\right)\approx 0.04288241\,\mathrm{rad}}\]

\item Conservar coseno en el equilibrio normal.

\[\fitmath{\mathrm{Vexact}=\sqrt{\frac{2\cdot 455\cdot 9.80665\cdot \cos\left(\delta\right)}{1.05\cdot 16.2\cdot 0.78}}}\]

\[\fitmath{\mathrm{Vexact}=\sqrt{\frac{2\cdot 455\cdot 9.80665\cdot \cos\left((0.04288241309)\right)}{1.05\cdot 16.2\cdot 0.78}}\approx 25.9228\,\mathrm{m/s}}\]

\item Proyectar la velocidad sobre la vertical.

\[\fitmath{\mathrm{sexact}=\mathrm{Vexact}\cdot \sin\left(\delta\right)}\]

\[\fitmath{\mathrm{sexact}=(25.92279697)\cdot \sin\left((0.04288241309)\right)\approx 1.111291\,\mathrm{m/s}}\]

\item Aplicar solo en esta ruta el equilibrio aproximado L=W.

\[\fitmath{\mathrm{Vapprox}=\sqrt{\frac{2\cdot 455\cdot 9.80665}{1.05\cdot 16.2\cdot 0.78}}\approx 25.93472\,\mathrm{m/s}}\]

\item Sustituir seno por tangente en la aproximación.

\[\fitmath{\mathrm{sapprox}=\frac{\mathrm{Vapprox}\cdot C_D}{0.78}}\]

\[\fitmath{\mathrm{sapprox}=\frac{(25.93472073)\cdot (0.0334688)}{0.78}\approx 1.112826\,\mathrm{m/s}}\]

\item Cuantificar el error en descenso.

\[\fitmath{\mathrm{relative}=\frac{\mathrm{sapprox}}{\mathrm{sexact}}-1}\]

\[\fitmath{\mathrm{relative}=\frac{(1.112825617)}{(1.111291422)}-1\approx 0.001380551\,}\]

\end{enumerate}





\apartado{10. Resultados con unidades}{Los resultados pertenecen al modelo y al intervalo de datos declarados.}

\begin{center}\begin{tabular}{lr}\toprule Magnitud&Resultado\\\midrule

TAS exacta & $25.9228\,\mathrm{m\,s^{-1}}$\\

Descenso exacto & $1.111291\,\mathrm{m\,s^{-1}}$\\

Descenso aproximado & $1.112826\,\mathrm{m\,s^{-1}}$\\

Error relativo de descenso & $0.001380551\,$\\

\bottomrule\end{tabular}\end{center}

\apartado{11. Interpretación}{La aproximación es buena para una eficiencia elevada, pero la igualdad exacta y la aproximada no son la misma relación.}

\apartado{12. Comprobación independiente y alcance}{Comprobar que el cociente de velocidades es sqrt(cos(delta)) y evaluar energía DV/W en el caso exacto.}

\apartado{13. Limitaciones}{Un error pequeño de trigonometría no compensa incertidumbre grande de la polar o falta de trim.}

\par\smallskip\noindent{\footnotesize\textit{Datos y solución originales. Las sustituciones registradas se reproducen con SymPy y Studium; eso no establece validación experimental ni revisión experta.}}
